\section{微積分の基本定理}\label{節:基本定理第一}
ここまで,定積分をリーマン和の極限として定義し,和から分かる性質を調べた.次に,右端を変数にして積分そのものを関数として見る.\,$f$ が $[a,b]$ で連続なとき,
\[
 A(x)=\int_a^x f(t)\dd t\qquad(a\le x\le b)
\]
とおく.\,$x$ が動くと,上端までに足し合わせる範囲が変わる.\,$f\ge0$ なら面積を表すが,負の値を含む場合にも符号付きの量として同じ式を使う.

\begin{figure}[htbp]
\centering
\begin{tikzpicture}[x=2.5cm,y=1.5cm,>=Stealth]
 \fill[blue!12] (0,0)--plot[domain=0:2,samples=60] (\x,{1+0.2*\x*\x})--(2,0)--cycle;
 \fill[red!18] (2,0)--plot[domain=2:2.5,samples=35] (\x,{1+0.2*\x*\x})--(2.5,0)--cycle;
 \draw[->] (-0.15,0)--(3.15,0) node[right]{$t$};
 \draw[->] (0,-0.1)--(0,3.15) node[above]{$y$};
 \draw[very thick,domain=0:3,samples=80] plot (\x,{1+0.2*\x*\x});
 \draw[dashed] (2,0)--(2,1.8);
 \draw[dashed] (2.5,0)--(2.5,2.25);
 \node[below] at (0,0) {$a$};
 \node[below] at (2,0) {$x$};
 \node[below] at (2.5,0) {$x+h$};
 \node[blue!70!black] at (1,0.55) {$A(x)$};
 \node[red!70!black] at (2.25,0.8) {増分};
 \node at (2.7,3.05) {$y=f(t)$};
\end{tikzpicture}
\caption{上端を $h$ だけ動かしたときの積分の増分は,幅 $h$ の区間の積分になる.}
\label{図:面積関数増分}
\end{figure}

図\ref{図:面積関数増分}は $h>0$ で $f$ が正の場合を描いている.幅が小さいとき,増分は高さ $f(x)$,幅 $h$ の長方形に近い.これを符号や $h$ の向きに依存しない形で示す.

\begin{theorem}{微積分の基本定理：積分を微分する}\label{定理:基本定理第一}
$f$ が $[a,b]$ で連続なら,\,$A(x)=\int_a^x f(t)\dd t$ は $[a,b]$ で連続であり,内点で微分可能で,
\[
 A'(x)=f(x)\qquad(a<x<b)
\]
を満たす.端点では,区間の内側からの片側微分で同じ式が成り立つ.
\end{theorem}

まず,連続関数 $f$ は閉区間で有界だから,\,$|f(t)|\le M$ となる $M$ がある.積分の評価と区間の加法性により,任意の $u,v\in[a,b]$ について
\[
 |A(v)-A(u)|=\left|\int_u^v f(t)\dd t\right|
 \le M|v-u|
\]
となる.よって $A$ は連続である.

微分を調べるため,\,$x,x+h\in[a,b]$,\,$h\ne0$ とする.区間の加法性から
\begin{align*}
 \frac{A(x+h)-A(x)}h-f(x)
 &=\frac1h\int_x^{x+h}f(t)\dd t-f(x)\\
 &=\frac1h\int_x^{x+h}\{f(t)-f(x)\}\dd t.
\end{align*}
ここで,定数関数 $f(x)$ の $x$ から $x+h$ までの積分が $hf(x)$ になることを使った.

任意の $\varepsilon>0$ に対し,\,$f$ の $x$ での連続性から,\,$|t-x|<\delta$ なら $|f(t)-f(x)|<\varepsilon$ となる $\delta>0$ を選べる.\,$0<|h|<\delta$ のとき,\,$x$ と $x+h$ の間のすべての $t$ がこの条件を満たす.積分区間の長さは $|h|$ だから,向きがどちらであっても
\[
 \left|\frac1h\int_x^{x+h}\{f(t)-f(x)\}\dd t\right|
 \le\frac1{|h|}\,\varepsilon|h|=\varepsilon.
\]
したがって差商は $f(x)$ に近づき,\,$A'(x)=f(x)$ を得る.図から予想した「積分の増加率が高さになる」という関係が,負の高さや $h<0$ の場合も含めて証明できた.



\subsection{和と差}
微分して元の関数になることを,数列の和でも考えてみよう.\,$S_n=\sum_{k=1}^n a_k$なら,
\[
 S_n-S_{n-1}=a_n
\]
である.和の隣り合う二つを引くと,最後に加えた一項が残る.一方,積分では
\[
 A(x+h)-A(x)=\int_x^{x+h}f(t)\dd t
\]
となり,新しく加えた短い区間が残る.\,$h$で割って極限を取れば,\,$A'(x)=f(x)$である.数列では一項分の差を取り,関数では短い区間の変化率を考えるという違いはあるが,いずれも,足したものからその場所で加わった量を取り出していると読める.

例えば速度$v(t)=2t$なら,原始関数は$t^2$なので,\,$\int_0^T2t\dd t=T^2$である.導入では長方形の和と三角形の面積から求めた値が,今度は微分の公式から得られた.一般の速度でも,原始関数が見つかれば同じ計算ができる.

%
